\section{Математические модели в биологии и биоинформатике}
\label{sec:models-in-biology-and-bioinformatics}

Математические модели в биологии описывают изменение численности
популяций, взаимодействие видов, эпидемические процессы, реакционно-
диффузионные системы и другие динамические явления. В биоинформатике
математические методы применяются к последовательностям ДНК, РНК и белков,
к филогенетическим данным и большим биологическим наборам данных.

\subsection{Модели роста популяции}

Простейшая модель Мальтуса предполагает постоянную удельную скорость роста:
\[
\boxed{\frac{dN}{dt}=rN},
\qquad
N(t)=N_0e^{rt}.
\]
Она пригодна лишь пока ограниченность ресурсов несущественна.

Модель Ферхюльста учитывает насыщение:
\[
\boxed{
\frac{dN}{dt}=rN\left(1-\frac{N}{K}\right),
}
\]
где $K$ --- ёмкость среды. Стационарные состояния:
\[
N_1^*=0,
\qquad
N_2^*=K.
\]
При $r>0$ состояние $N=0$ неустойчиво, а $N=K$ асимптотически устойчиво.
Таким образом, даже простая нелинейность качественно меняет долгосрочный
прогноз.

\subsection{Взаимодействие видов: модель Лотки--Вольтерры}

Классическая модель ``хищник--жертва'' имеет вид
\[
\boxed{
\begin{aligned}
\dot x&=ax-bxy,\\
\dot y&=-cy+dxy,
\end{aligned}
}
\qquad a,b,c,d>0,
\]
где $x$ --- численность жертв, $y$ --- хищников.

Кроме тривиального равновесия существует положительное стационарное
состояние
\[
\boxed{
(x^*,y^*)=\left(\frac{c}{d},\frac{a}{b}\right).
}
\]
Линеаризация системы около стационарного состояния позволяет исследовать
его устойчивость. В идеальной классической модели возникают нейтральные
колебания; в более реалистичных моделях добавляют внутривидовую
конкуренцию, насыщение хищника и другие эффекты.

\subsection{Эпидемиологическая модель SIR}

Пусть $S(t)$, $I(t)$ и $R(t)$ --- числа восприимчивых, инфицированных и
удалённых из эпидемического процесса особей, а
\[
N=S+I+R=\text{const}.
\]
Одна из стандартных форм модели:
\[
\boxed{
\begin{aligned}
\dot S&=-\beta\frac{SI}{N},\\
\dot I&=\beta\frac{SI}{N}-\gamma I,\\
\dot R&=\gamma I.
\end{aligned}
}
\]
Базовое репродуктивное число
\[
\boxed{R_0=\frac{\beta}{\gamma}}
\]
показывает среднее число вторичных заражений в полностью восприимчивой
популяции. В начале эпидемии
\[
\dot I>0
\quad\Longleftrightarrow\quad
R_0\frac{S_0}{N}>1.
\]
При $S_0\approx N$ это даёт известный порог $R_0>1$.

\subsection{Пространственные модели}

Если существен пространственный перенос, к кинетике добавляют диффузию.
Например, для вектора концентраций $u(x,t)$:
\[
\frac{\partial u}{\partial t}
=D\Delta u+F(u).
\]
Такие \textbf{реакционно-диффузионные системы} используются для описания
распространения биологических популяций, морфогенеза, пространственных
эпидемий и химико-биологических волн.

\subsection{Биоинформатика: выравнивание последовательностей}

Пусть даны две последовательности
\[
A=a_1\ldots a_n,
\qquad
B=b_1\ldots b_m.
\]
Задача глобального выравнивания состоит в поиске расположения совпадений,
замен и пропусков, максимизирующего суммарный балл. Она решается методом
динамического программирования. Для штрафа $d>0$ за пропуск и функции
сходства $s(a_i,b_j)$ рекуррентное соотношение Нидлмана--Вунша имеет вид
\[
\boxed{
F_{ij}=\max\left\{
\begin{array}{l}
F_{i-1,j-1}+s(a_i,b_j),\\
F_{i-1,j}-d,\\
F_{i,j-1}-d
\end{array}
\right\}.
}
\]
Граничные условия:
\[
F_{i0}=-id,
\qquad
F_{0j}=-jd.
\]
После заполнения таблицы оптимальное выравнивание восстанавливается
обратным проходом. Сложность классического алгоритма составляет
$O(nm)$ по времени. Для локального выравнивания алгоритм Смита--Уотермана
дополнительно допускает обнуление отрицательного значения.

\subsection{Другие модели биоинформатики}

В зависимости от задачи применяют вероятностные и статистические модели:
скрытые марковские модели для аннотации последовательностей, модели
эволюции для филогенетики, регрессионные и машинно-обучающие методы для
анализа высокоразмерных биологических данных. Важно, что математическая
модель всегда связана с конкретной биологической гипотезой и должна
проверяться на независимых данных.

\subsection{Что важно сказать на экзамене}

Удобно привести три разных типа моделей: динамическую модель роста
(логистическое уравнение), систему взаимодействующих популяций
Лотки--Вольтерры и SIR. Для биоинформатики хорошим конкретным примером
является выравнивание последовательностей и рекуррентная формула
динамического программирования. Желательно уметь объяснить физический или
биологический смысл параметров и стационарных состояний.

\newpage
